\section{Classical propositional rejection}
\label{annex:p02-classical-rejection}

\begin{annexitem}{annex:p02.crf.substitution}{thm:p02-crf-substitution}{The valuation pullback is a structural induction}
\begin{proof}
On a variable \(x\), the equality is the definition
\(w(\sigma(x))=v(x)\). If it holds for \(B\), then
\(w(\sigma(\neg B))=1-w(\sigma(B))=1-v(B)=v(\neg B)\).
If it holds for \(B,C\), Boolean implication has value zero exactly
when its antecedent has value one and consequent value zero.
Therefore \(w(\sigma(B)\to\sigma(C))\) and \(v(B\to C)\)
have the same value. These cases exhaust the grammar.
If \(w(\sigma(A))=0\), the constructed \(v\) consequently
refutes \(A\). If \(A\) is valid, every such \(v\) gives one,
so every \(w\) makes \(\sigma(A)\) true. Positive completeness
then supplies a closed proof of \(\sigma(A)\). This last argument
uses positive completeness, not substitution into open hypotheses.
\end{proof}
\end{annexitem}

\begin{annexitem}{annex:p02.crf.soundness}{thm:p02-crf-soundness}{Rejection soundness}
\begin{proof}
Induct on the finite rejection derivation, keeping the invariant that
every negative line has a countervaluation. The antiaxiom \(p\)
has one: assign \(p\) value zero and assign any values to the
remaining variables.
For rejection modus tollens, a preceding negative line \(B\) has
a countervaluation \(w\). The accompanying closed proof of
\(A\to B\) makes that implication true at every valuation by
positive soundness. At \(w\), its consequent is zero, so its
antecedent must be zero. Thus the same \(w\) refutes \(A\).
For reverse substitution, a countervaluation of \(\sigma(A)\)
pulls back to one of \(A\) by the previous result. These are all
negative-production cases. Positive lines require no negative witness
and are valid by their checked closed proofs. The invariant therefore
holds at the last line. A valid formula cannot have such a witness;
positive soundness also excludes simultaneous assertion and rejection
of the same formula.
\end{proof}
\end{annexitem}

\begin{annexitem}{annex:p02.crf.completeness}{thm:p02-crf-completeness}{Rejection completeness}
\begin{proof}
Let \(v(A)=0\). For every valuation \(w\), the formulas
\(T=(p\to p)\) and \(F=\neg(p\to p)\) have values one
and zero. Consequently \(w(\sigma_v(x))=v(x)\) for each variable
of \(A\). The substitution lemma yields
\(w(\sigma_v(A))=v(A)=0\) for every \(w\). Simultaneous
replacement makes this argument valid even when \(p\) occurs in
\(A\) or in an image of the substitution. No iterative substitution
equation is being solved.
An implication with antecedent always zero is always true, so
\(\sigma_v(A)\to p\) is a tautology. The empty-premise positive
adequacy theorem supplies a finite closed Hilbert proof of it.
Attach that proof as a positive certificate. The antiaxiom
\(\dashv p\), followed by rejection modus tollens, gives
\(\dashv\sigma_v(A)\); reverse substitution then gives
\(\dashv A\). Conversely, a derivable rejection has a
countervaluation by the soundness theorem. This proves both
directions. The source's formula \(p\land\neg p\) has value
zero at both Boolean values of \(p\), just like \(F\), so the
stated truth-function bridge between the two constant presentations
also holds.

For the worked formula \(q\to r\), the false assignment produces
\(S=(T\to F)\). Here is an explicit derivation of \(p\) under
the sole temporary hypothesis \(S\). A1--A3 are the actual schemas
on \cref{sheet:hilbert-cpl}; \(F=\neg T\).
\begin{longtable}{@{}r p{0.64\textwidth} p{0.22\textwidth}@{}}
\caption{Hilbert derivation of \(p\) under the temporary hypothesis \(S\)}\\
\toprule Line & Formula & Justification\\\midrule
\endfirsthead
\caption[]{Hilbert derivation of \(p\) under the temporary hypothesis \(S\) (continued)}\\
\toprule Line & Formula & Justification\\\midrule
\endhead
\midrule\multicolumn{3}{r}{\emph{continued on next page}}\\
\endfoot
\endlastfoot
1 & \(p\to((p\to p)\to p)\) & A1\\
2 & \(p\to(p\to p)\) & A1\\
3 & \((p\to((p\to p)\to p))\to((p\to(p\to p))\to(p\to p))\) & A2\\
4 & \((p\to(p\to p))\to(p\to p)\) & MP, 1 and 3\\
5 & \(T=p\to p\) & MP, 2 and 4\\
6 & \(S=T\to F\) & Temporary hypothesis\\
7 & \(F\) & MP, 5 and 6\\
8 & \(F\to(\neg p\to F)\) & A1\\
9 & \(\neg p\to F\) & MP, 7 and 8\\
10 & \(T\to(\neg p\to T)\) & A1\\
11 & \(\neg p\to T\) & MP, 5 and 10\\
12 & \((\neg p\to\neg T)\to((\neg p\to T)\to p)\) & A3\\
13 & \((\neg p\to T)\to p\) & MP, 9 and 12\\
14 & \(p\) & MP, 11 and 13\\
\bottomrule
\end{longtable}
The constructive deduction theorem
\cref{thm:p02-hilbert-deduction}, proved in
\cref{annex:p2.calc.2}, converts this sequence into a closed
proof of \(S\to p\). Only that discharged proof is supplied to
rejection modus tollens. The positive hypothesis in line 6 does not
survive as an open assumption of the rejection certificate.
\end{proof}
\end{annexitem}

\begin{annexitem}{annex:p02.crf.polarity}{thm:p02-crf-polarity}{Quantifiers, contingent formulas, and illegal rule directions}
\begin{proof}
The soundness/completeness pair for rejection gives
\(\dashv A\) exactly when some valuation makes \(A\) zero.
Positive soundness/completeness gives \(\vdash\neg A\)
exactly when every valuation makes \(\neg A\) one, equivalently
when every valuation makes \(A\) zero. A formula therefore has
one of three Boolean truth functions: always one, always zero, or
one at some valuations and zero at others. The first case gives
assertion of \(A\) and excludes its rejection. The second gives
rejection of \(A\) and assertion of \(\neg A\). In the third,
a zero valuation of \(A\) rejects \(A\), and a one valuation
of \(A\) rejects \(\neg A\). Neither formula is valid, so
neither has a positive proof. For \(A=p\), valuations with
\(p=0\) and \(p=1\) are the two distinct witnesses. They are
not a single valuation making both formulas false.

For forward negative substitution, the permitted antiaxiom is
\(\dashv p\). Substitution \(p\mapsto p\to p\) would give
rejection of the tautology \(p\to p\) if applied forward to
that line. Hence such a rule violates rejection soundness.
For the hypothesis-dependent positive-premise variant, take the
hypothesis \(q\). A1 gives \(q\to(T\to q)\), and modus
ponens gives \(T\to q\) under that hypothesis. The genuine
calculus rejects \(q\): apply reverse substitution to the
antiaxiom using \(q\mapsto p\). Using the open proof of
\(T\to q\) with this rejection would reject the tautology
\(T\). It is therefore essential that the positive premise
be a closed theorem, as required by the actual rule.
\end{proof}
\end{annexitem}

\begin{annexitem}{annex:p02.crf.procedure}{thm:p02-crf-procedure}{Finite semantic search and effective proof enumeration}
\begin{proof}
A finite formula contains finitely many variables, found by a syntax
traversal. For \(n\) distinct variables there are \(2^n\)
Boolean assignments. Recursive evaluation terminates on each one,
so the truth-table phase terminates and decides whether the input
is valid. Assignments outside that finite variable set have no
effect on its value, by induction on the formula grammar.

Finite formulas and finite formula sequences over the effectively
encoded countable alphabet can be enumerated. For each candidate
sequence, axiom-schema membership is decidable by matching its finite
syntax against A1--A3 with repeated metavariables required to match
identically. Each proposed modus-ponens step is checked by its two
earlier indices and exact formulas. Thus closed proofs are
effectively enumerable without an undecided validity oracle.
In the valid branch, positive completeness ensures a finite proof
of \(A\), so the enumeration eventually finds one. In the
invalid branch, the constructed substitution is finite on the
variables of \(A\) and the completeness proof above establishes
that \(\sigma_v(A)\to p\) has a finite closed proof. The
second enumeration therefore also terminates. The appended
negative steps are finite and their substitutions are checkable
by recursive syntax traversal.

Positive soundness establishes the meaning of an accepted positive
certificate; rejection soundness establishes the meaning of an
accepted negative one. The termination argument uses existence
of a finite proof, not a polynomial bound or a claim that brute-force
proof enumeration is a practical solver. The \(q\to r\) example
already has a true valuation as well as the false one, demonstrating
why a negative output does not establish universal falsity.
\end{proof}
\end{annexitem}

\begin{annexitem}{annex:p02.crf.disjunction}{thm:p02-crf-disjunction}{The logical condition for a disjunctive rejection}
\begin{proof}
Suppose the displayed disjunction property holds and \(A,B\notin L\).
If \(A\lor B\in L\), that property would put at least one of \(A,B\)
in \(L\), a contradiction. Thus \(A\lor B\notin L\).

Conversely, suppose the rejection inference preserves nonmembership.
If \(A\lor B\in L\) but neither disjunct belonged to \(L\), its two
negative premises would hold while its negative conclusion failed.
Therefore at least one disjunct belongs to \(L\). This establishes the
equivalence, using classical reasoning about membership in \(L\);
it makes no assertion that the object logic \(L\) is classical.

For classical tautologies take \(A=p\) and \(B=\neg p\).
Each has a falsifying valuation, so neither is in \(L\), whereas
\(p\lor\neg p\) is true at both Boolean values of \(p\).
The inference therefore fails. In contrast, if one valuation makes
both disjuncts false, the Boolean truth condition for disjunction
makes their disjunction false at that same valuation.
\end{proof}
\end{annexitem}

\begin{annexitem}{annex:p02.crf.join}{thm:p02-crf-join}{Joining the supplied assignments}
\begin{proof}
If the finite assignments agree wherever their domains overlap, their
union assigns exactly one value to each variable in the union of those
domains. Extending it by value zero outside that domain gives a total
Boolean valuation. The value of a formula depends only on its occurring
variables: the variable case is immediate, and each connective computes
its value from the values of its immediate subformulas. Structural
induction therefore shows that the extended valuation preserves each
formula's value under its original assignment. All listed formulas
remain false. Repeated application of the Boolean disjunction clause
makes their finite disjunction false.

Conversely, a common extension must have the value specified by each
assignment on every overlap. Conflicting values on any shared variable
therefore exclude a common extension. This proves the claimed equivalence.

To see why failed joining is not a satisfiability verdict, consider
\(A=p\land q\) and \(B=\neg p\land q\). One countervaluation of \(A\)
assigns \(p=0,q=1\); one countervaluation of \(B\) assigns \(p=1,q=1\).
These supplied assignments conflict on \(p\). Nevertheless \(p=0,q=0\)
makes both formulas false, so alternative compatible countervaluations
exist. The procedure tests the supplied assignments, not that broader
existence question.
\end{proof}
\end{annexitem}
